% ============================================================
%  QCE26 Workshop Paper
%  IEEE Quantum Week 2026 — Toronto, Sep 13–18, 2026
%
%  Requirements (source: qce.quantum.ieee.org/2026/call-for-workshops/):
%    - \documentclass[10pt,conference]{IEEEtran}  (mandatory)
%    - NO compsoc / compsocconf option
%    - Workshop papers: 4 pages TOTAL, references included
%    - Submit via EasyChair: https://easychair.org/conferences/?conf=qce26
% ============================================================

\documentclass[10pt,conference]{IEEEtran}

% ── Core packages ────────────────────────────────────────────
\usepackage[T1]{fontenc}
\usepackage[utf8]{inputenc}
\usepackage{microtype}          % better line breaking
\usepackage{cite}               % IEEE-style compressed citations [1]–[3]
\usepackage{amsmath,amssymb,amsthm}
\usepackage{graphicx}
\usepackage{booktabs}           % \toprule / \midrule / \bottomrule
\usepackage{algorithm}
\usepackage{algpseudocode}      % algorithmicx pseudocode
\usepackage{xcolor}
\usepackage{hyperref}
\hypersetup{
    colorlinks = true,
    linkcolor  = blue!70!black,
    citecolor  = blue!70!black,
    urlcolor   = blue!70!black,
}

% ── Theorem environments ─────────────────────────────────────
\newtheorem{theorem}{Theorem}
\newtheorem{lemma}[theorem]{Lemma}
\newtheorem{proposition}[theorem]{Proposition}
\newtheorem{corollary}[theorem]{Corollary}
\theoremstyle{definition}
\newtheorem{definition}{Definition}
\newtheorem{problem}{Problem}
\theoremstyle{remark}
\newtheorem{remark}{Remark}

% ── Qubit mapping / routing macros ───────────────────────────
\newcommand{\G}{\mathcal{G}}            % coupling graph
\newcommand{\A}{\mathcal{A}}            % architecture graph
\newcommand{\C}{\mathcal{C}}            % circuit
\newcommand{\Q}{\mathcal{Q}}            % qubit set
\newcommand{\M}{\mathcal{M}}            % mapping / layout
\newcommand{\SWAP}{\mathrm{SWAP}}
\newcommand{\CNOT}{\mathrm{CNOT}}
\newcommand{\depth}{\mathrm{depth}}
\newcommand{\dist}{\mathrm{dist}}       % graph distance
\newcommand{\NP}{\mathsf{NP}}
\newcommand{\todo}[1]{\textcolor{red}{\textbf{[TODO: #1]}}}

% Placeholder box for figures whose final artwork is dropped in later.
\newcommand{\figslot}[2][0.95\linewidth]{%
    \setlength{\fboxsep}{0pt}%
    \fbox{\parbox[c][#2][c]{#1}{\centering\footnotesize\itshape
        \textcolor{gray}{[ figure placeholder — \detokenize{#1} wide ]}}}%
}

% ── Metadata ─────────────────────────────────────────────────
\title{%
    Qubit Mapping and Routing for Fault-Tolerant Compilation:\\
    \large A Gaussian Potential-Field Placement Heuristic%
}

\author{%
    \IEEEauthorblockN{%
        Alessandro Ruzza\IEEEauthorrefmark{1},
        Federico Quartieri\IEEEauthorrefmark{1},
        Daniele Salvi\IEEEauthorrefmark{1}
    }
    \IEEEauthorblockA{
        \IEEEauthorrefmark{1}
        Politecnico di Milano, Milan, Italy\\
        \texttt{\{alessandro.ruzza, federico.quartieri, daniele.salvi\}@polimi.it}
    }
}

% ============================================================
\begin{document}

\maketitle

% ── Abstract ─────────────────────────────────────────────────
\begin{abstract}
Executing a logical quantum circuit fault-tolerantly on a surface-code
processor reduces to placing logical qubits on a two-dimensional lattice and
scheduling lattice-surgery operations that consume magic states for
non-Clifford gates. Unlike the near-term (NISQ) setting, the cost is dominated
not by $\SWAP{}$ insertion but by the time-multiplexed routing corridors that
must be kept free so that interacting patches, and patches in need of a magic
state, can be connected on demand. We present a \emph{Gaussian potential-field}
heuristic for the initial-placement (mapping) sub-problem of fault-tolerant
compilation. Each candidate location on the lattice is scored by a
superposition of two-dimensional Gaussian influence fields: attractive fields
pull T-gate--heavy qubits toward magic states and tightly interacting qubits
toward their already-placed partners, while repulsive (inverted) fields spread
qubits apart and preserve routing space. Placement is greedy in qubit-priority
order and is filtered by a \emph{safe-passage} connectivity test that
guarantees every committed patch can still reach its interaction partners and a
magic state. We implement the method in an open-source C++20 compiler that
ingests OpenQASM and emits a routed lattice-surgery schedule, and we evaluate it
against a rule-based magic-aware placement, a fast randomized-restart baseline,
and the external WISQ compiler on a suite of benchmark circuits.
\todo{Insert one or two headline numbers once tables are filled (e.g.\ mean
routing-step reduction vs.\ baseline; compile time; circuits where WISQ
times out but our method succeeds).}
\end{abstract}

% IEEE keywords
\begin{IEEEkeywords}
qubit mapping, qubit routing, fault-tolerant compilation, surface code,
lattice surgery, magic states, quantum circuit compilation, heuristic placement
\end{IEEEkeywords}

% ============================================================
\section{Introduction}
\label{sec:intro}

Fault-tolerant quantum computation encodes each logical qubit in a patch of a
quantum error-correcting code---most commonly the surface code, laid out on a
two-dimensional grid of physical qubits~\cite{fowler2012surface}. Logical
two-qubit gates are realised by \emph{lattice surgery}: a strip of ancilla
patches between two operands is temporarily merged and split to perform a
joint measurement~\cite{horsman2012lattice,litinski2019game}. Non-Clifford
gates (the $T$ gate in the Clifford$+T$ universal set) cannot be applied
transversally and are instead teleported by consuming a distilled \emph{magic
state}, which must be delivered from a dedicated magic-state factory through a
free corridor of the lattice~\cite{bravyi2005magic}.

Compiling a logical circuit to such hardware therefore has two coupled
sub-problems. The \emph{mapping} problem assigns logical qubits to patches on
the lattice, and the \emph{routing} problem time-schedules the lattice-surgery
operations, reserving disjoint ancilla corridors at each step. This mirrors the
classical NISQ mapping-and-routing pipeline of compilers such as
SABRE~\cite{li2019sabre} and the Qiskit transpiler~\cite{qiskit}, but the cost model is fundamentally
different: there are no $\SWAP{}$ gates, connectivity is reconfigured every step,
and the scarce resource is free \emph{space} for corridors and magic-state
access rather than edges of a fixed coupling graph. Both sub-problems are
$\NP$-hard in their classical analogues~\cite{siraichi2018qubit}, and the
fault-tolerant variants inherit this hardness, motivating fast heuristics.

A good initial placement is decisive: if T-heavy qubits sit far from magic
states, or strongly interacting qubits are scattered, every routing step pays
for it in serialised lattice surgery. Existing fault-tolerant toolchains such as
WISQ~\cite{wisq} and lattice-surgery schedulers~\cite{beverland2022surface}
typically treat placement and routing jointly with search or solver-based
methods that scale poorly. We instead ask whether a cheap, geometry-aware
\emph{placement} heuristic can produce layouts that route well under a simple
scheduler.

\smallskip
\noindent\textbf{Contributions.}
\begin{itemize}
    \item We formulate fault-tolerant initial placement as the maximisation of
    a \emph{Gaussian potential field} over the lattice: a per-qubit score
    surface built by superposing attractive and repulsive 2D Gaussians sourced
    by magic states, mapped interaction partners, and already-placed qubits
    (Section~\ref{sec:approach}).
    \item We couple the placement with a \emph{safe-passage} connectivity
    filter that rejects any placement which would block a committed patch from
    reaching its CNOT partners or a magic state, guaranteeing a routable layout.
    \item We give two parameterisations---a discrete \emph{coarse} variant and a
    continuous \emph{fine} variant that interpolates field strength from
    per-qubit $T$-count and per-pair CNOT statistics.
    \item We provide an open-source C++20 implementation (OpenQASM
    front-end, lattice-surgery router, benchmark harness) and evaluate it
    against rule-based magic-aware and randomized-restart baselines and the
    WISQ compiler (Section~\ref{sec:eval}).
\end{itemize}

The rest of the paper is organised as follows.
Section~\ref{sec:background} fixes the architecture and circuit model and
states the placement problem. Section~\ref{sec:approach} describes the Gaussian
field method. Section~\ref{sec:eval} reports the experimental evaluation, with
space reserved for the Gaussian-mapping result tables and figures.
Sections~\ref{sec:related} and~\ref{sec:conclusion} discuss related work and
conclude.

% ============================================================
\section{Background and Problem Formulation}
\label{sec:background}

\subsection{Architecture and Circuit Model}

We target a rectangular $W\times H$ grid of surface-code patches. We model it as
an \emph{architecture graph} $\A = (V, E)$ whose vertices $V$ are tiles, each
with integer coordinates $(x,y)$, and whose edges $E$ connect orthogonally
adjacent tiles (the corridors along which lattice surgery can merge patches). A
fixed subset $V_{\mathrm{m}} \subset V$ of tiles host \emph{magic states};
their number and placement (a central cluster or a boundary row) are
configuration parameters. The remaining tiles are available for logical-qubit
patches or as transient routing corridors.

A circuit $\C$ over logical qubits $\Q_L$ is taken in the Clifford$+T$ basis.
From $\C$ we extract, per logical qubit $q$, its $T$-count $t(q)$ (number of
non-Clifford gates, i.e.\ magic-state demand) and, for each pair $(q,q')$, the
two-qubit interaction count $c(q,q')$. We write $\bar t,\sigma_t$ for the mean
and standard deviation of the $T$-count across qubits, and $\bar c,\sigma_c$
for the corresponding CNOT-count statistics; these summary statistics drive the
field weights in Section~\ref{sec:approach}.

\subsection{Fault-Tolerant Routing Model}

A layout $\M:\Q_L\to V$ is injective: each logical qubit occupies one tile. The
router processes the circuit in layers and, at each step, must realise the
active two-qubit gates by reserving vertex-disjoint paths in $\A$ between
operand tiles, and must realise active $T$ gates by reserving a path from the
operand tile to some free magic-state tile. Gates whose corridors collide are
serialised into later steps. The primary cost we minimise is the number of
\emph{routing steps} (the lattice-surgery time-depth); we also report the
average gate parallelism (gates per step) and the fraction of gates that could
not be routed at first exposure.

\subsection{Safe Passage}

Because corridors are reclaimed each step, the layout must never \emph{trap} a
patch. We say a placement satisfies \emph{safe passage} if, after committing it,
(i) every pair of already-placed CNOT partners is still connected by some free
path, (ii) every placed patch with remaining $T$ demand can still reach a magic
state, and (iii) at least one magic-state tile keeps a free neighbouring entry
for the not-yet-placed $T$-consuming qubits. Safe passage is the
fault-tolerant analogue of the connectivity constraint that ordinary routing
takes for granted, and it is what makes a purely score-driven placement
admissible.

\begin{definition}[Qubit Mapping]
    Given circuit $\C$ and architecture $\A$, a \emph{mapping}
    $\M : \Q_L \to V\setminus V_{\mathrm m}$ is an injective assignment of
    logical qubits to non-magic tiles.
\end{definition}

\begin{problem}[Routable Placement]
    Find a mapping $\M$ that (a) satisfies safe passage and (b) minimises the
    expected number of lattice-surgery routing steps required to execute $\C$
    on $\A$.
\end{problem}

% ============================================================
\section{Proposed Approach: Gaussian Field Placement}
\label{sec:approach}

The core idea is to convert the soft, competing objectives of fault-tolerant
placement---``be near magic states if T-heavy'', ``cluster with frequent
partners'', ``leave room to route''---into an additive scalar field on the
lattice, then place each qubit at the field's safe maximum.

\subsection{Influence Fields}
\label{sec:fields}

A single \emph{influence field} is an axis-aligned 2D Gaussian centred at a tile
$(\mu_x,\mu_y)$,
\begin{equation}
g(x,y) \;=\; w\,\exp\!\Big(\!-\tfrac{(x-\mu_x)^2}{2\sigma_x^2}
                              -\tfrac{(y-\mu_y)^2}{2\sigma_y^2}\Big),
\label{eq:gauss}
\end{equation}
with non-negative weight $w$. A field is either \emph{attractive}, contributing
$g(x,y)$, or \emph{repulsive} (``inverted''), contributing the valley
$\max\!\big(w-g(x,y),\,0\big)$, which is small at the centre and rises to $w$
far away. The widths are set from a single interpretable parameter, the
\emph{confidence} $\kappa\in(0,1)$: with $R$ the half-extent of the grid along
an axis,
\begin{equation}
\sigma \;=\; \frac{R}{\sqrt{-2\ln(1-\kappa)}},
\label{eq:sigma}
\end{equation}
so that the field retains a fraction $(1-\kappa)$ of its peak at distance $R$
from the centre; larger $\kappa$ yields broader, longer-range fields. A
\texttt{size} multiplier scales the supported footprint independently.

For the qubit $q$ currently being placed, the total score at a free tile $v$ is
\begin{equation}
S(v) \;=\;
   \underbrace{\sum_{m\in V_{\mathrm m}}\! g^{\mathrm{mag}}_m(v)}_{\text{magic-state pull/push}}
 + \underbrace{\sum_{p\in P(q)}\! g^{\mathrm{cnot}}_p(v)}_{\text{partner clustering}}
 + \underbrace{\sum_{u\in \M}\! g^{\mathrm{rep}}_u(v)}_{\text{spread-out}}
 + \underbrace{b(v)}_{\text{centering}},
\label{eq:score}
\end{equation}
optionally plus a constant border bonus. The four families are:

\paragraph{Magic-state fields}
one Gaussian per magic tile, whose sign and weight encode $q$'s $T$ demand:
T-heavy qubits make them attractive (pull toward magic states), T-light qubits
make them repulsive (push away, freeing magic neighbourhoods for qubits that
need them), and average qubits zero them out.

\paragraph{CNOT (partner) fields}
one attractive Gaussian centred on the already-mapped tile of each frequent
interaction partner $p\in P(q)$, so that strongly coupled qubits coalesce and
their future surgeries traverse short corridors.

\paragraph{Repulsion (mapped) fields}
one repulsive Gaussian centred on every already-placed qubit, discouraging
overcrowding and preserving routing space---the field-based surrogate for
congestion avoidance.

\paragraph{Baseline field}
a single attractive Gaussian at the grid centre that keeps the overall layout
compact.

Figure~\ref{fig:fields} illustrates the four families and their superposition
into the score surface $S$.

% ---- FIGURE: influence-field decomposition (single column) ----
\begin{figure}[t]
    \centering
    \figslot[0.95\textwidth]{1.5in}
    \caption{Decomposition of the Gaussian placement score for one qubit:
    magic-state, CNOT-partner, repulsion, and baseline fields (left) and their
    superposition $S$ over the lattice (right). The chosen tile is the safe
    maximum of $S$. \todo{Replace with rendered field maps from
    \texttt{visualization/gaussian\_frames/}.}}
    \label{fig:fields}
\end{figure}

\subsection{Weight Assignment: Coarse and Fine}

The two variants differ only in how the per-field weights $w$ and the
attractive/repulsive sign are chosen from circuit statistics.

\emph{Coarse} uses discrete decisions. If $t(q)$ exceeds $q$'s maximum CNOT
count, the magic fields are made strongly attractive (weight \textsc{magic-high});
otherwise placement is partner-driven: attractive CNOT fields (weight
\textsc{cnot-high}) are added at mapped partners, and the magic fields are
attractive when $t(q)>\bar t+\sigma_t$, repulsive when $t(q)<\bar t-\sigma_t$,
and disabled in between.

\emph{Fine} replaces the thresholds with linear interpolation. For the magic
fields,
\begin{equation*}
w \;=\;
\begin{cases}
\textsc{magic-high}, & \hspace{-1cm} |t(q)-\bar t|>\sigma_t,\\[2pt]
\mathrm{lerp}\big(\textsc{magic-low},\textsc{magic-high},\tfrac{|t(q)-\bar t|}{\sigma_t}\big), & \hspace{-0.2cm} \text{otherwise},
\end{cases}
\end{equation*}
attractive when $t(q)>\bar t$ and repulsive when $t(q)<\bar t$ (zero at the
mean). The CNOT fields are weighted analogously from the per-pair count
$c(q,p)$ relative to $\bar c,\sigma_c$. Fine thus reacts smoothly to how far a
qubit deviates from the circuit's average resource profile.

\subsection{Placement Loop}

Logical qubits are placed greedily in priority order: a max-heap ranks qubits by
their interaction importance---the sum of a qubit's $T$-count and its largest
single-partner CNOT count---so that the most resource-hungry qubits are committed
first, while their relations are still free to shape. For each popped qubit, the
score field $S$ of \eqref{eq:score} is evaluated over all free tiles; tiles are
sorted by descending $S$ (stable, with a deterministic tie-break) and the first
one passing the safe-passage test of Section~\ref{sec:background} is committed.
A new repulsion field is then added at the chosen tile. The procedure is
summarised in Algorithm~\ref{alg:main}; Figure~\ref{fig:frames} shows snapshots
of a layout forming.

\begin{algorithm}[t]
\caption{Gaussian field placement}
\label{alg:main}
\begin{algorithmic}[1]
    \Require Circuit $\C$, architecture $\A$ with magic tiles $V_{\mathrm m}$
    \Ensure Safe-passage mapping $\M$
    \State build priority heap of $\Q_L$ by interaction importance
    \State $\M \gets \emptyset$;\; add baseline field; add magic fields
    \While{heap not empty}
        \State $q \gets$ pop highest-priority unmapped qubit
        \State set magic/CNOT field weights from $q$'s $T$- and CNOT-statistics
                \Comment{coarse or fine}
        \State $S(v)\gets$ superpose fields over each free tile $v$ \Comment{Eq.~\eqref{eq:score}}
        \For{$v$ in tiles sorted by $S(v)$ descending}
            \If{\Call{SafePassage}{$v,q,\M$}}
                \State $\M(q)\gets v$;\; add repulsion field at $v$;\; \textbf{break}
            \EndIf
        \EndFor
        \If{no tile committed} \State \textbf{fail} (no routable placement) \EndIf
    \EndWhile
    \State \Return $\M$
\end{algorithmic}
\end{algorithm}

% ---- FIGURE: layout-formation frames (full width, top) ----
\begin{figure*}[t]
    \centering
    \figslot[0.95\textwidth]{1.7in}
    \caption{Formation of a Gaussian-field layout for a representative
    benchmark (e.g.\ \texttt{qft\_20}). Each frame shows the current score
    surface $S$ and the committed patches; magic tiles in the centre, repulsion
    fields visibly spacing qubits apart while T-heavy qubits cluster around the
    magic cluster. \todo{Replace with stills exported by the
    \texttt{make-video} target.}}
    \label{fig:frames}
\end{figure*}

\begin{remark}
Field evaluation is $O(|V|\cdot(|V_{\mathrm m}|+|P(q)|+|\M|))$ per qubit; with
the lattice sized as $\Theta(|\Q_L|+|V_{\mathrm m}|)$ tiles, placement is
low-order polynomial. In practice the per-candidate safe-passage connectivity
checks dominate placement time. \todo{Insert measured placement/compile times
from the benchmark CSVs.}
\end{remark} 

% ============================================================
\section{Evaluation}
\label{sec:eval}

\subsection{Experimental Setup}

\textbf{Implementation.} The method is implemented in an open-source C++20
compiler that parses OpenQASM, generates the rectangular architecture and
magic-state placement, runs the mapping of Section~\ref{sec:approach}, and
routes the result with a lattice-surgery router (naive or congestion-aware
shortest-path corridors, with a dedicated $T$-gate-to-magic routing strategy).
The lattice dimensions are sized automatically from the qubit count, the
interaction degree, and the magic-state budget.

\textbf{Benchmarks.} We use \todo{$N$} circuits spanning arithmetic
(\texttt{adder}, \texttt{multiplier}, \texttt{square\_root}), algorithmic
(\texttt{qft}, \texttt{qpe}, \texttt{bwt}, \texttt{hhl}, \texttt{factor247}),
and structured/state-prep (\texttt{ghz}, \texttt{wstate}, \texttt{ising},
\texttt{qaoa}) families, with logical qubit counts from \todo{min} to
\todo{max}.

\textbf{Baselines and metrics.} We compare the coarse and fine Gaussian variants
against three baselines that share the same architecture generator and router:
(i) \emph{magic-aware}, a rule-based greedy placement that sends each qubit to a
magic-state neighbourhood or beside its heaviest CNOT partner according to
threshold rules on its $T$- and CNOT-counts; (ii) \emph{random}, a fast
placement that scatters qubits over a sparse, magic-avoiding sublattice whose
enforced spacing makes every layout trivially routable; and (iii) the external
WISQ compiler~\cite{wisq}. The random baseline is the natural speed reference:
because each restart is cheap---no field evaluation and no per-candidate
connectivity test---it is run for many randomized repetitions and the
best-routing layout is kept, whereas the Gaussian and magic-aware placements are
deterministic and run once. Random buys this speed with grid area: lacking any
circuit-structure information, it is allotted a larger, sparser lattice, while
the Gaussian and magic-aware methods pack onto a roughly one-third smaller grid
and recover routability through the safe-passage test. The primary metric is the
number of routing steps; we also report compile time, average parallelism, and
the non-routed-layer percentage.

\subsection{Results}

\paragraph{Routing quality}
Table~\ref{tab:results} reports routing steps and compile time across the
benchmark suite for both Gaussian variants and the three baselines.
Figure~\ref{fig:scatter} plots per-circuit routing steps of the Gaussian method
against WISQ; points below the diagonal favour our method, and the rightmost
column of Table~\ref{tab:results} flags circuits where WISQ fails or times out
while the Gaussian method completes.

% ---- TABLE: main comparison (skeleton, ample room) ----
\begin{table}[t]
    \centering
    \caption{Routing steps and compile time per benchmark. ``G-coarse'' and
    ``G-fine'' are the Gaussian variants; ``MA'' is the rule-based magic-aware
    baseline; ``Rand'' is the best of $k$ randomized restarts; WISQ is the
    external compiler (\,---\,$=$ failure/timeout).
    \todo{Fill from \texttt{best\_params\_benchmark\_runs.csv} and
    \texttt{wisq\_docker} CSVs.}}
    \label{tab:results}
    \setlength{\tabcolsep}{4pt}
    \footnotesize
    \begin{tabular}{l r rrrrr r}
        \toprule
        \textbf{Circuit} & \textbf{$|\Q_L|$}
            & \multicolumn{5}{c}{\textbf{Routing steps}}
            & \textbf{Time} \\
        \cmidrule(lr){3-7}
            & & \textbf{G-coarse} & \textbf{G-fine} & \textbf{MA} & \textbf{Rand} & \textbf{WISQ}
            & \textbf{(s)} \\
        \midrule
        \texttt{adder\_n28}        & -- & -- & -- & -- & -- & -- & -- \\
        \texttt{qft\_20}           & -- & -- & -- & -- & -- & -- & -- \\
        \texttt{multiplier\_n45}   & -- & -- & -- & -- & -- & -- & -- \\
        \texttt{square\_root\_n45} & -- & -- & -- & -- & -- & -- & -- \\
        \texttt{bwt\_n21}          & -- & -- & -- & -- & -- & -- & -- \\
        \texttt{hhl\_n10}          & -- & -- & -- & -- & -- & -- & -- \\
        \texttt{factor247\_n15}    & -- & -- & -- & -- & -- & -- & -- \\
        ~~$\vdots$                 &    &    &    &    &    &    &    \\
        \midrule
        \textbf{Mean / geomean}    & -- & -- & -- & -- & -- & -- & -- \\
        \bottomrule
    \end{tabular}
\end{table}

% ---- FIGURE: per-circuit scatter vs WISQ (single column) ----
\begin{figure}[t]
    \centering
    \figslot{1.6in}
    \caption{Per-circuit routing steps: Gaussian field placement vs.\
    WISQ~\cite{wisq}. Points below the diagonal indicate fewer routing steps for
    our method. \todo{Generate from merged benchmark CSVs.}}
    \label{fig:scatter}
\end{figure}

\paragraph{Parameter sensitivity}
The fields are governed by a small set of parameters: the confidence $\kappa$ of
\eqref{eq:sigma} (field width), the size multiplier (footprint), and the
high/low magic and CNOT weights. Table~\ref{tab:sweep} summarises a sweep over
the two most influential knobs, $\kappa$ and size; in our preliminary tuning
$\kappa$ chiefly controls effective field radius and interacts with the
safe-passage margin, with a broad, robust optimum.
\todo{Confirm best-found setting and headline sensitivity once filled.}

% ---- TABLE: parameter sensitivity (skeleton) ----
\begin{table}[t]
    \centering
    \caption{Sensitivity of mean routing steps to confidence $\kappa$ and
    Gaussian size multiplier (fine variant). \todo{Fill from
    \texttt{size\_confidence\_tuning} CSVs.}}
    \label{tab:sweep}
    \setlength{\tabcolsep}{6pt}
    \footnotesize
    \begin{tabular}{l cccc}
        \toprule
        \textbf{size $\backslash\ \kappa$} & $0.90$ & $0.95$ & $0.99$ & $0.999$ \\
        \midrule
        $0.5$ & -- & -- & -- & -- \\
        $0.7$ & -- & -- & -- & -- \\
        $1.0$ & -- & -- & -- & -- \\
        \bottomrule
    \end{tabular}
\end{table}

\paragraph{Discussion}
\todo{2--4 sentences once numbers are in: (i) coarse vs.\ fine trade-off;
(ii) where the field method beats / loses to magic-aware and WISQ and why
(magic-state proximity vs.\ corridor crowding); (iii) the value of
circuit-structure information --- how many randomized restarts the random
baseline needs to match field placement, if ever, and at what grid-area cost;
(iv) scalability --- circuits on which WISQ times out but field placement
finishes in seconds.}

% ============================================================
\section{Related Work}
\label{sec:related}

\textbf{NISQ mapping and routing.} Heuristic layout-and-$\SWAP{}$ methods such as
SABRE~\cite{li2019sabre}, the routing passes of
\mbox{t$|$ket$\rangle$}~\cite{sivarajah2020tket} and Qiskit~\cite{qiskit}, and
exact/$\NP$-hardness treatments~\cite{siraichi2018qubit} target a fixed coupling
graph and minimise $\SWAP{}$ count or depth. Our cost model has no $\SWAP{}$s and a
per-step reconfigurable topology, so these objectives do not transfer directly,
though the greedy, lookahead-driven spirit does.

\textbf{Fault-tolerant compilation.} Lattice surgery~\cite{horsman2012lattice}
and its surface-code resource analyses~\cite{litinski2019game,
beverland2022surface} establish the routing primitives and cost units we
optimise; magic-state distillation and consumption~\cite{bravyi2005magic} make
T-gate placement the central tension our magic-state fields address. WISQ~\cite{wisq}
is a recent end-to-end fault-tolerant compiler that we use as an external
baseline. Relative to solver- and search-heavy schedulers, our contribution is
a cheap, interpretable \emph{placement} prior---the superposed Gaussian field---
that produces routable layouts for a simple downstream router.

\textbf{Potential-field methods.} Artificial potential fields are classic in
robot motion planning; we adapt the attractive/repulsive-field idea to discrete
lattice placement, where magic states and partners act as attractors and placed
patches as repellers, under a hard connectivity (safe-passage) constraint.

% ============================================================
\section{Conclusion}
\label{sec:conclusion}

We presented a Gaussian potential-field heuristic for the initial-placement
sub-problem of fault-tolerant compilation. By superposing attractive
magic-state and partner fields with repulsive spread-out fields and selecting
the safe maximum, the method turns the competing objectives of T-gate locality,
interaction clustering, and corridor preservation into a single, fast scoring
pass, while a safe-passage filter guarantees routable layouts. Implemented in an
open-source OpenQASM-to-lattice-surgery compiler, it is evaluated against a
rule-based baseline and the WISQ compiler.
\todo{State the headline empirical result once tables are filled.}
Limitations include the greedy, single-pass nature of placement and reliance on
hand-tuned field weights. A natural next step is to learn or jointly optimise
the field weights and to feed the field gradient back into the router so that
placement and scheduling co-adapt.

% ── References ───────────────────────────────────────────────
\bibliographystyle{IEEEtran}
\bibliography{bib}

\end{document}
